%% Springer Nature LaTeX template (sn-jnl) — single-file submission.
%% Prepared for: Empirical Economics (Springer).
%% Reference style: author-year (sn-mathphys-ay) via natbib (\citet/\citep).
%% Single-blind submission: author identities included on the title page.

\documentclass[pdflatex,sn-mathphys-ay]{sn-jnl}

%%%% Standard packages
\usepackage[T1]{fontenc}% needed for accented names (e.g. Zi\k{e}ba) in the bibliography
\usepackage{graphicx}%
\usepackage{multirow}%
\usepackage{amsmath,amssymb,amsfonts}%
\usepackage{amsthm}%
\usepackage{mathrsfs}%
\usepackage[title]{appendix}%
\usepackage{xcolor}%
\usepackage{textcomp}%
\usepackage{manyfoot}%
\usepackage{booktabs}%

\raggedbottom

\begin{document}

\title[Identity, Not Count: Model Risk in Parsimonious Credit Overlays]{Identity, Not Count: Indicator Selection as the First-Order Model Risk in Radically Reduced Credit Overlays}

%%=============================================================%%
%% Single-blind submission: author identities included.
%%=============================================================%%

\author[1]{\fnm{Subin} \sur{Ji}}\email{subinji33@ecredible.co.kr}

\author[2]{\fnm{Hyunjeong} \sur{Kang}}\email{hyunjkang@kookmin.ac.kr}

\author*[3]{\fnm{Munil} \sur{Yang}}\email{miyang@ips.or.kr}

\affil[1]{\orgname{ECREDIBLE Co., Ltd.}, \orgaddress{\city{Seoul}, \country{Republic of Korea}}.
ORCID: \url{https://orcid.org/0009-0003-8904-6346}}

\affil[2]{\orgname{Kookmin University}, \orgaddress{\city{Seoul}, \country{Republic of Korea}}.
ORCID: \url{https://orcid.org/0009-0007-9181-4649}}

\affil[3]{\orgname{Institute for Industrial Policy Studies}, \orgaddress{\city{Seoul}, \country{Republic of Korea}}.
ORCID: \url{https://orcid.org/0009-0004-1493-6691}}

%%==================================%%
%% Abstract (unstructured; no citations/equations per SN guidance)       %%
%%==================================%%
\abstract{How much of the specification risk in a reduced credit model comes from how \emph{many}
indicators it uses, and how much from \emph{which}? We study a parsimonious overlay---two financial
indicators laid over a grade built from the full ratio set---on public data with a reproducible,
out-of-fold design, and show its value is governed by which indicators enter, not by the fact that there
are two. Across 66 indicator pairs the incremental likelihood ratio spans two orders of
magnitude: informative pairs dominate while uninformative ones are indistinguishable from zero, and the
ordering survives Bonferroni and false-discovery control, though every pair adds exactly two regressors.
The operative property is informativeness, not a textbook category: a permutation test does not find the
solvency category to outrank the liquidity/activity category. The identity-over-count pattern replicates
across a bankruptcy panel, a rating cross-section, and four independent consumer-credit datasets. By contrast the overlay's incremental
\emph{ranking} power is modest and fragile---significant on the pooled panel but only near default by
horizon, and not robust to a continuous base probability of default---so we claim no uniform
discrimination gain. A rating-recovery analysis shows two indicators are least adequate at the rating
extremes and most adequate in the middle band (bootstrap $95\%$ CI excludes zero); a downside-asymmetry
pattern appears in the point estimates but is not established under the correct contrast, so we report it
only as suggestive. The contribution is a reproducible demonstration that, in a reduced-form overlay,
indicator selection is the first-order source of specification risk.}

\keywords{Credit risk, Default prediction, Model risk, Feature selection, Specification risk, Reproducibility}

\pacs[JEL Classification]{C52, C53, G21, G24, G28, G32}

\maketitle

%%=====================================================%%
\section{Introduction}\label{sec:intro}

When a credit model is radically reduced---refreshed between full reviews using only a handful of
cheap indicators---where does its specification risk come from: from how \emph{many} indicators it uses,
or from \emph{which}? This is the question we study, through a deliberately parsimonious overlay of two
financial indicators laid over a grade built from the full ratio set. The credit-scoring literature has
long treated the choice and number of predictors as a first-order design problem rather than a detail
\citep{lessmann2015benchmarking,dumitrescu2022machine}, and a feature-selection literature shows that a
small, well-chosen set can preserve most of a model's accuracy. We push that logic to its extreme---two
indicators---and ask, within the tightest possible reduction, whether it is the count or the identity of
the indicators that governs how much information the overlay carries. Our answer is that identity
dominates: two indicators that carry solvency information behave entirely differently from two that do
not, even though both add exactly two regressors, and this distinction---not the number two---is what a
validation process for a reduced overlay should scrutinize. We test this with out-of-fold statistics on
fully public data and corroborate it across a bankruptcy panel, a rating cross-section, and four
independent consumer-credit datasets.

A natural setting for this question is the common practice of assessing credit on two clocks. A full
assessment, based on the complete set of financial statements, is set once a year; between annual
revisions, the grade may be refreshed using only a small number of indicators that are cheap to
collect. This practice is common in supply-chain and SME settings, where a large buyer or a lender
re-scores many suppliers within the year but can afford a full statement-based review only annually.
It produces two grades for the same firm---an annual grade and an interim grade---that need not
agree, and the discrepancy is frequently treated as a problem to be reconciled rather than as a risk
signal to be used. The object of this paper is the default-risk content of that reduced overlay:
used as a low-cost device to monitor firms for deterioration between full reviews, does it actually
measure default risk that the annual grade misses, and can that content be validated on public data?

We argue that the two-clock structure is better understood through the lens of rating philosophy.
The annual, full-statement grade is effectively a \emph{through-the-cycle} (TTC) assessment that
smooths transitory fluctuations, whereas the interim update built from a few indicators is a
low-cost \emph{point-in-time} (PIT) overlay \citep{altman2004stability,altman2006pit,loffler2004anatomy};
Figure~\ref{fig:framework} illustrates the resulting two-clock structure.
Framed this way, the operational question---why do the two grades differ?---becomes an empirical
one about credit-risk measurement: does a deliberately parsimonious PIT overlay carry default-risk
information beyond the TTC grade, and if so, for which firms? This is the question we study.

\begin{figure}[t]
\centering
\includegraphics[width=0.92\textwidth]{figures/fig_framework_ttc_pit}
\caption{Conceptual framework. An annual full-statement grade is a through-the-cycle (TTC)
assessment; an interim update built from two indicators is a point-in-time (PIT) overlay. The paper
asks whether the parsimonious PIT overlay adds information beyond the TTC grade.}
\label{fig:framework}
\end{figure}

Three features make the question non-trivial. First, the interim overlay uses only \emph{two}
indicators, so whatever information it adds must survive extreme parsimony. Second, in practice the
overlay is applied selectively---typically to the middle of the rating spectrum, while the
investment-grade and near-default extremes are left to the annual review---and it is not obvious
whether this selection is a prudent default or an avoidable loss of information. Third, the value of
a two-indicator signal may be asymmetric: the credit literature consistently finds that markets
respond to deterioration far more than to improvement \citep{holthausen1986,dichev2001}.

Studying interim re-rating directly would require proprietary agency data with an observed mid-year
revision for each firm. Absent such data, we develop a fully reproducible design on two public datasets
that \emph{reconstructs} the two-clock structure rather than observing it: we treat the operational
practice as a methodological hypothesis and test a feature-parsimony analogue of it, being explicit
throughout that the temporal interpretation is exactly that---an interpretation (we return to this
construct-validity limitation in Section~\ref{sec:discussion}). The first dataset is a corporate bankruptcy panel with a
realized default outcome \citep{zieba2016}; the second is a cross-section of corporate credit
ratings with agency letter grades. Because one dataset has a default outcome but no ratings and the
other has ratings but no default outcome, the two analyses form a complementary
internal/external-validity pair rather than two independent exercises.

Our findings are as follows. The central result is that the overlay's value depends decisively on
\emph{which} two indicators are chosen rather than on their number: across 66 indicator pairs the
incremental likelihood ratio spans two orders of magnitude, the informative pairs dominate and the
uninformative ones are indistinguishable from zero, and the ordering survives false-discovery control.
The ordering reflects informativeness rather than a textbook category---a permutation test does not find
the solvency category to outrank the liquidity/activity category---so in a reduced-form overlay the
selection of indicators, not their count, is the first-order source of specification risk. The two
supporting findings frame this one honestly. First, the overlay's incremental ranking power over the
base grade is modest and fragile: significant on the pooled panel but, split by horizon, only near
default, and not robust to a continuous base probability of default rather than a discretized grade, so
we do not claim a uniform discrimination gain. Second, in the rating data (Study~2) a nonparametric
recovery analysis shows that two indicators recover roughly half of the rating information in the full
ratio set, with the loss significantly larger at the rating extremes than in the middle band, backed by
a bootstrap confidence interval---an independent confirmation, on a different dataset and outcome, of
the middle-band pattern. A downside pattern, in which deterioration appears more informative than
improvement, is present in the point estimates but is not statistically established once the correct
equal-slopes contrast is tested, and we report it only as a suggestive regularity.

These results speak to the measurement, validation, and governance of credit risk. To the TTC/PIT
literature \citep{altman2004stability,loffler2004anatomy,topp2010ttcpit} we add the question of whether a
\emph{parsimonious} PIT overlay measures default risk beyond the TTC grade, and an answer that depends on
the rating band. To the rating-transition and surveillance literature
\citep{lando2002transitions,guttler2007adjustment} we offer an incremental-information test that
distinguishes risk signal from noise in within-year grade changes. To the model-risk and
model-validation literature \citep{bcbs2005validation,sr2011modelrisk,blochlinger2006,danielsson2002} we
contribute a reproducible validation protocol for a reduced-form overlay and show that, within such an
overlay, indicator identity rather than indicator count is a first-order source of specification risk. And to the SME and
supply-chain credit literature \citep{altmansabato2007,berger2006} we provide a public-data method for
evaluating interim re-rating policies without proprietary data. Throughout, the contribution is as much
methodological as substantive: we show how the credit-risk content of a parsimonious overlay on a
statement-based grade can be measured, tested, decomposed, and stress-tested for model risk with
out-of-fold discrimination statistics on fully public data---the kind of transparent, reproducible
validation that supervisory guidance on model risk management calls for \citep{sr2011modelrisk}.

The remainder of the paper is organized as follows. Section~\ref{sec:literature} reviews the related
literature. Section~\ref{sec:data} describes the two datasets. Section~\ref{sec:method} sets out the
research design, hypotheses, and econometric models. Section~\ref{sec:results} reports the main
results, Section~\ref{sec:robustness} the robustness analysis, Section~\ref{sec:discussion} the
implications, and Section~\ref{sec:conclusion} concludes.

%%=====================================================%%
\section{Related literature}\label{sec:literature}

Our study connects six strands of the credit-risk literature: the philosophy of
through-the-cycle versus point-in-time ratings; the dynamics of rating transitions and
surveillance; default prediction and rating replication from financial ratios; the role of
parsimony in credit models; the asymmetry of credit information; and the model risk and
validation of quantitative credit models. We also draw on the SME and supply-chain credit
setting that motivates interim re-rating, and on the methodology for measuring and comparing
model discrimination.

\subsection{Through-the-cycle versus point-in-time ratings}
Rating agencies deliberately smooth their grades over the business cycle. \citet{altman2004stability}
show that agencies achieve stability through a prudent-migration policy and a change
threshold, trading timeliness for fewer rating reversals, and \citet{altman2006pit} quantify
how far such through-the-cycle (TTC) grades lag a point-in-time (PIT) benchmark.
\citet{loffler2004anatomy} formalizes the TTC filter and the conditions under which it helps or
harms default prediction, and later asks whether agencies can genuinely separate trend from
cycle \citep{loffler2013lookthrough}. \citet{topp2010ttcpit} contrast the two philosophies and the
problem of mapping between them, while the agency-side statement of the accuracy--stability
trade-off is given by \citet{cantor2006tradeoff}. This literature frames our central object: an
annual full-statement grade is effectively a TTC assessment, and a low-cost interim update is a
PIT overlay. What it does not address is whether a \emph{deliberately parsimonious} PIT overlay,
built from only a few indicators, adds information beyond the TTC grade.

\subsection{Rating transitions, surveillance, and the noise-versus-information question}
Whether a within-year rating change reflects information or noise is an old question.
\citet{lando2002transitions} estimate migration intensities in continuous time and document
rating drift, and \citet{fledelius2004nonparametric} show transition behavior varies with
time-in-rating. \citet{nickell2000stability} and \citet{bangia2002migration} establish that
transition matrices are not stationary but depend on industry and the business cycle, so that an
observed grade change partly reflects cyclical conditions. \citet{guttler2007adjustment} document
that agencies accelerate downgrades as default approaches, evidence that rating changes do carry
information near distress. \citet{carey2001parameterizing} caution that default probabilities
inferred from rating grades can be biased and unstable. Our paper reframes the ``grade change =
risk'' concern operationally: rather than asking whether a mid-year change is noise, we test
whether a two-indicator interim signal is \emph{incrementally predictive} of default once the
annual grade is controlled for.

\subsection{Default prediction and rating replication from financial ratios}
The prediction of corporate default from financial ratios begins with the discriminant-analysis
Z-score of \citet{altman1968}, the logit O-score of \citet{ohlson1980}, and the probit model of
\citet{zmijewski1984}; \citet{shumway2001} shows that a discrete-time hazard formulation corrects
the bias of single-period models. Machine-learning extensions---random forests and gradient
boosting \citep{barboza2017,jones2017}, and the ensemble-of-boosted-trees model of
\citet{zieba2016} that accompanies the Polish bankruptcy data we use---improve discrimination
further. A parallel literature replicates \emph{ratings} rather than defaults from fundamentals,
from the ordered-probit models of \citet{kaplan1979} and \citet{blume1998} to machine-learning
rating predictors \citep{huang2004,shin1999,kim2012,golbayani2020}. We build on both: Study~1
constructs a base grade and tests incremental default prediction, while Study~2 measures how much
rating information two indicators recover.

\subsection{Parsimony and feature selection in credit models}
How many and which financial ratios a credit model needs is itself a research question.
\citet{handhenley1997} review the construction of parsimonious scorecards, and a feature-selection
literature \citep{tsai2009,dujardin2010,liang2015,hajek2013,laborda2021}, surveyed by
\citet{lin2012}, shows that a small, well-chosen set of variables can preserve most of a model's
accuracy and that the \emph{choice} of variables can matter as much as the classifier. This
motivates our treatment of the interim overlay as a two-indicator problem and our finding that the
identity of the two indicators is decisive.

\subsection{Asymmetry of credit information}
A robust empirical regularity is that credit markets react more to bad news than to good news.
\citet{holthausen1986} and \citet{hand1992} find significant negative stock responses to
downgrades but weak or insignificant responses to upgrades; \citet{goh1993} identify which
downgrades convey bad news; and \citet{dichev2001} document negative long-run returns after
downgrades with no reliable drift after upgrades. The asymmetry extends to CDS and bond markets
\citep{hull2004,norden2004} and strengthens after Regulation FD \citep{jorion2005}. This evidence
anchors our hypothesis that a two-indicator signal predicts default asymmetrically: deterioration
should matter more than improvement.

\subsection{SME and supply-chain credit; validation methodology}
Interim re-rating is most consequential for SME suppliers assessed by large buyers.
\citet{altmansabato2007} show that SME-specific default models outperform generic corporate models,
\citet{altman2010nonfinancial} add non-financial information, and \citet{berger2006} map the
lending technologies---including trade-credit and credit-scoring lending---used in this segment.
The trade-credit and supply-chain-finance literature \citep{petersen1997,mian1992,giannetti2011,wuttke2013}
describes the buyer--supplier financing relationships in which repeated, low-cost credit
assessment arises. Finally, we rely on established tools for measuring and comparing
discrimination: the AUC and its interpretation \citep{hanley1982}, the nonparametric test for
correlated AUCs of \citet{delong1988}, the accuracy-ratio/CAP framework and tests of
\citet{engelmann2003}, model-validation guidance \citep{stein2007,krahnen2001}, and standard
logistic-regression inference \citep{hosmer2013}.

\subsection{Model risk and the validation of reduced-form credit overlays}
Whether a reduced-form overlay measures credit risk reliably is a model-risk and model-validation
question. The credit-risk measurement tradition ranges from the structural default model of
\citet{merton1974} to reduced-form intensity models \citep{jarrowturnbull1995,duffiesingleton1999} and
the rating-system architectures that operationalize them inside institutions \citep{crouhy2001prototype}.
Against this backdrop, supervisory and methodological work stresses that any such model must be
\emph{validated}: the \citet{bcbs2005validation} studies set out discrimination and calibration
standards for internal rating systems, \citet{sr2011modelrisk} frames model risk management---independent
validation, testing of limitations, and governance---as a supervisory expectation, and
\citet{sobehart2000benchmarking} and \citet{blochlinger2006} develop benchmarking and economic-value
criteria for default models. A parallel caution holds that quantitative risk models can mislead when
their specification is taken for granted \citep{danielsson2002,derman1996modelrisk}. Our out-of-fold,
leakage-free protocol is a validation exercise in exactly this sense, and our indicator- and
estimator-robustness analyses (Section~\ref{sec:robustness}) quantify the specification component of model
risk in a radically reduced overlay---showing that \emph{which} indicators are chosen, not how many, is
the dominant source of model risk.

\subsection{Contribution}
Relative to this literature we make five contributions. First, we formalize a widely used but
little-studied practice---a parsimonious PIT overlay on a TTC grade---as a \emph{credit-risk measurement}
question and test it directly with out-of-fold discrimination statistics. Second, and centrally, we show
with a 66-pair placebo distribution, surviving false-discovery control and replicated across five
datasets, that \emph{which} two indicators are used---their informativeness---is a first-order design
choice and the dominant source of specification-related model risk, not a detail, while the mere count
of indicators is not. Third, we document that the overlay's recovery of rating information is
concentrated in the middle band and weakest at the extremes, backed by a bootstrap confidence interval,
which rationalizes the common practice of excluding rating extremes from interim re-rating. Fourth, we treat the exercise as a model-validation protocol---out-of-fold, leakage-free, and
robust across estimators---aligned with supervisory model-risk guidance \citep{bcbs2005validation,sr2011modelrisk}.
Fifth, we provide a fully reproducible, public-data methodology for studying and governing interim
re-rating without proprietary agency data.

%%=====================================================%%
\section{Data}\label{sec:data}

We use two public datasets. One has a realized default outcome but no ratings; the other has agency
ratings but no default outcome. Their complementary strengths motivate the two-study design.

\subsection{Study 1: corporate bankruptcy panel}
The first dataset is the Polish companies bankruptcy data of \citet{zieba2016}, distributed through
the UCI Machine Learning Repository. It comprises five files indexed by forecast horizon: the
$n$-th file contains financial ratios observed in year $n$ of the window and a binary indicator of
bankruptcy within the subsequent $6-n$ years, so the fifth file corresponds to a one-year-ahead
default horizon. Pooling the five horizons gives 43{,}405 firm-year observations with 2{,}091
bankruptcies (an overall rate of 4.82\%); the rate rises monotonically from 3.86\% at the longest
horizon to 6.94\% at the one-year horizon (Table~\ref{tab:polish}), a gradient we exploit as a
pseudo-time dimension in the robustness analysis. Each observation carries 64 financial ratios.
Following the solvency emphasis of \citet{altman1968}, we take the two proxy indicators to be a
profitability ratio (net profit / total assets) and a leverage ratio (total liabilities / total
assets); Figure~\ref{fig:proxy} shows both separate bankrupt from non-bankrupt firms in the expected
direction. Missingness is low overall ($\approx 1.3$--$1.9\%$) except for one ratio with 39--45\%
missing values; crucially, this heavily missing ratio is \emph{not} one of the two proxy indicators,
both of which are almost fully observed. We median-impute missing values and add a missingness
indicator for the heavily missing ratio, so that the base grade does not confound imputation with
signal; results are unchanged if that ratio is dropped from the base grade entirely.

\begin{table}[t]
\centering
\caption{Study 1 sample by forecast horizon (Polish bankruptcy data).}
\label{tab:polish}
\begin{tabular}{@{}lrrr@{}}
\toprule
Horizon file & Observations & Bankruptcies & Rate (\%) \\
\midrule
1 (5-year)   & 7{,}027  & 271 & 3.86 \\
2            & 10{,}173 & 400 & 3.93 \\
3            & 10{,}503 & 495 & 4.71 \\
4            & 9{,}792  & 515 & 5.26 \\
5 (1-year)   & 5{,}910  & 410 & 6.94 \\
\midrule
Pooled       & 43{,}405 & 2{,}091 & 4.82 \\
\botrule
\end{tabular}
\end{table}

\begin{figure}[t]
\centering
\includegraphics[width=0.78\textwidth]{figures/fig02_polish_proxy_by_class}
\caption{Distribution of the two proxy indicators by bankruptcy status (one-year horizon).}
\label{fig:proxy}
\end{figure}

\subsection{Study 2: corporate credit ratings}
The second dataset is a cross-section of 2{,}029 corporate credit ratings compiled from public
sources (OpenML dataset \#43344), spanning 2005--2016 and twelve sectors, with 25 financial ratios
per firm and no missing values. Ratings are issued by five agencies---Standard \& Poor's (the
plurality), Egan-Jones, Moody's, Fitch, and DBRS---so we include agency fixed effects throughout and
report a Standard \& Poor's--only specification as a robustness check. The letter grades range from
AAA to D; the middle band (BBB, BB, B) accounts for 72.1\% of observations and the extremes for the
remainder (Figure~\ref{fig:ratingdist}), reproducing the concentration that makes middle-band interim
re-rating the operationally relevant case. For the ordinal analysis we collapse the three rarest
grades (CC, C, D; eight observations in total) into CCC, yielding a stable seven-grade scale. The two
proxy indicators analogous to Study~1 are the debt ratio (leverage) and return on assets
(profitability).

\begin{figure}[t]
\centering
\includegraphics[width=0.72\textwidth]{figures/fig03_sp_rating_distribution}
\caption{Rating distribution; middle-band grades (BBB/BB/B) are shaded darker.}
\label{fig:ratingdist}
\end{figure}

%%=====================================================%%
\section{Research design and methodology}\label{sec:method}

Figure~\ref{fig:design} summarizes the design: two public datasets with complementary strengths provide
complementary evidence on the information value of a parsimonious interim overlay. The datasets share no
firms and differ in country, period, and outcome, so they are not two estimates of one quantity; rather,
Study~1 speaks to incremental default prediction and Study~2 to rating-information recovery, and we read
agreement between them as mutual corroboration rather than replication.

\begin{figure}[t]
\centering
\includegraphics[width=0.9\textwidth]{figures/fig_design_two_studies}
\caption{Research design. The two datasets form a complementary validity pair---one has a default
outcome but no ratings, the other ratings but no default outcome---and jointly test H1--H5.}
\label{fig:design}
\end{figure}

\subsection{Reconstructing a parsimony benchmark on public data}
We reconstruct a feature-parsimony benchmark on each dataset: a strong, information-rich base grade
built from the full ratio set, against which a deliberately two-indicator overlay is tested. The
two-clock interpretation---an annual full grade refreshed by a low-cost interim signal---motivates the
construction but is not observed in the data, and we return to this construct-validity limitation in
Section~\ref{sec:discussion}. In Study~1 the base grade is proxied by a
\emph{base grade} $G$ obtained by estimating a probability of default from the 62 ratios that
\emph{exclude} the two proxy indicators, via out-of-fold logistic regression, and discretizing the
fitted probability into seven quantile buckets (a scale analogous to the AAA--CCC range). Excluding
the two proxies from the base grade ensures that any incremental information attributed to the
two-indicator overlay is genuinely new. The overlay signal $S$ is the out-of-fold fitted default
probability from a logistic regression on the two proxies alone. The ``middle band'' is the set of
central buckets; the ``extremes'' are the rest. All out-of-fold quantities use five-fold stratified
cross-validation with a fixed seed. To avoid information leakage, the quantile bucket edges that define
the grade are computed within each training fold only and applied to the held-out fold; we report a
fully nested re-estimation in Section~\ref{sec:robustness} and confirm the main result survives it.
Figure~\ref{fig:protocol} maps each step of this construction. We emphasize that the two-clock,
annual-versus-interim reading is an \emph{interpretation} of a feature-parsimony design, not an
observed timeline; the identification here is cross-sectional, and the econometric object is the
incremental default-risk content of two indicators beyond a full-ratio grade. Section~\ref{sec:discussion}
discusses the resulting construct-validity limitation.

\begin{figure}[t]
\centering
\includegraphics[width=0.95\textwidth]{figures/fig_protocol_simulation}
\caption{Reconstruction pipeline (not a Monte Carlo simulation). The base grade is built from the
other ratios and the two-indicator overlay from the two indicators; the sample is restricted to the
middle band; and the final step tests whether the overlay adds information beyond the grade.}
\label{fig:protocol}
\end{figure}

\subsection{Hypotheses}
The design yields five testable hypotheses about the incremental default-risk content of a
two-indicator overlay:
\begin{itemize}
\item[H1.] Controlling for the base grade, a worse two-indicator signal predicts higher default risk.
\item[H2.] Adding the two-indicator signal to the base grade significantly improves discrimination (AUC / accuracy ratio).
\item[H3.] The incremental information is asymmetric: deterioration predicts default more strongly than improvement reduces it.
\item[H4.] The incremental value is larger in the middle band than at the extremes.
\item[H5.] In rating recovery, the information lost by using only two indicators is larger at the extremes than in the middle.
\end{itemize}

The asymmetry hypothesis (H3) is motivated by a simple intuition rather than a formal model, and we
treat it accordingly. The intuition is that, under a limited-liability payoff, the default-probability
function is locally steep in the deterioration region and flat once a firm is safely solvent, so a
reduced solvency signal might carry more information on the downside---in the same spirit as the market
evidence that bad credit news moves prices more than good \citep{holthausen1986,dichev2001}. We stress
that this is a motivating conjecture, not an identified mechanism: the sign pattern it predicts is also
produced mechanically whenever the signal is more dispersed on the downside, and as Section~\ref{sec:results}
shows we cannot statistically separate the two, so H3 is reported only as a suggestive direction. The
middle-band hypotheses (H4--H5), by contrast, are supported by evidence that does not rely on this
decomposition: at the investment-grade and near-default extremes two indicators are least able to add or
recover information, and the Study~2 recovery result that establishes this (H5) is backed by a bootstrap
confidence interval.

\subsection{Econometric models (Study 1)}
Default is modeled with a discrete-time logit including horizon fixed effects. The nested pair is
\begin{align}
\text{logit}\,\Pr(D_i=1) &= \alpha + \beta\,G_i + \boldsymbol{\delta}^{\top}X_i, \label{eq:m0}\\
\text{logit}\,\Pr(D_i=1) &= \alpha + \beta\,G_i + \gamma\,S_i + \boldsymbol{\delta}^{\top}X_i, \label{eq:m1}
\end{align}
where $X_i$ collects horizon dummies. H1 is tested by the sign and significance of $\gamma$ in
\eqref{eq:m1}; H2 by comparing the discrimination of \eqref{eq:m0} and \eqref{eq:m1} using the
\citet{delong1988} test on paired out-of-fold ROC curves and a likelihood-ratio test. For H3 we
decompose the overlay signal into its deviation from the grade-peer mean, $d_i = S_i - \bar{S}_{G(i)}$,
split into a deterioration part $d_i^{+}=\max(d_i,0)$ and an improvement part $d_i^{-}=\min(d_i,0)$:
\begin{equation}
\text{logit}\,\Pr(D_i=1) = \alpha + \beta\,G_i + \gamma^{+}d_i^{+} + \gamma^{-}d_i^{-} + \boldsymbol{\delta}^{\top}X_i, \label{eq:h3}
\end{equation}
and test $|\gamma^{+}| > |\gamma^{-}|$. For H4 we interact the signal with a middle-band indicator
$\text{Mid}_i$:
\begin{equation}
\text{logit}\,\Pr(D_i=1) = \alpha + \beta\,G_i + \gamma\,S_i + \eta\,(S_i\times\text{Mid}_i) + \theta\,\text{Mid}_i + \boldsymbol{\delta}^{\top}X_i, \label{eq:h4}
\end{equation}
with H4 corresponding to $\eta>0$.

\subsection{Rating recovery and information loss (Study 2)}
For Study~2 we measure how much rating information the two indicators recover relative to the full
ratio set. Predictive recovery is assessed nonparametrically with a random forest (300 trees, $\sqrt{p}$ features
per split, unrestricted depth, fixed seed) under five-fold cross-validation, scored by the quadratic
weighted kappa (QWK), exact accuracy, and within-one-grade accuracy, with QWK confidence intervals
obtained by 2{,}000-replication bootstrap. Statistical inference on the two proxies uses an ordered logit with agency fixed effects,
\begin{equation}
R_i^{*} = \boldsymbol{\theta}^{\top}Z_i + \varepsilon_i, \qquad R_i = r \iff \mu_{r-1} < R_i^{*} \le \mu_r, \label{eq:ordlogit}
\end{equation}
where $Z_i$ contains the two proxies and agency dummies. H5 is evaluated by comparing the
within-one-grade recovery of the two-indicator and full models \emph{by band}, and visualized with a
grade migration matrix. To benchmark the full ordinal model without the convergence problems of a
high-dimensional maximum-likelihood ordered logit, we also estimate an $L_2$-penalized
proportional-odds model (Appendix~\ref{app:po}); conclusions are invariant across estimators.

\subsection{Discrimination metrics}
Throughout we report the AUC and the accuracy ratio (Gini $=2\,\text{AUC}-1$), interpreting them per
\citet{hanley1982} and \citet{engelmann2003}, the Brier score for calibration, and QWK for ordinal
recovery. All discrimination comparisons are out-of-fold, and incremental comparisons use the
\citet{delong1988} and likelihood-ratio tests.

%%=====================================================%%
\section{Results}\label{sec:results}

\subsection{Study 1: incremental default prediction}
Table~\ref{tab:h1h2} reports the nested comparison. All figures in this paper come from a single
reproducible pipeline (fixed seeds; code released), so the pooled increment reported here matches the
pooled row of the horizon decomposition in Table~\ref{tab:horizon} up to rounding. Adding the two-indicator
signal raises the out-of-fold AUC from 0.729 to 0.739 and the accuracy ratio from 0.458 to 0.477. The
level improvement is small---unsurprisingly, since the base grade already uses 62 ratios---but on the
pooled panel it is significant: the \citet{delong1988} statistic is $z=7.77$ ($p=7.8\times10^{-15}$),
the likelihood-ratio statistic is 82.3 ($p<10^{-18}$), and the coefficient on $S$ in \eqref{eq:m1} is
positive and significant.
Figure~\ref{fig:roc} shows the two ROC curves. On the pooled panel, H1 and H2 are supported:
controlling for a grade built from all other ratios, the two indicators still carry default
information, and a fully nested, leakage-free re-estimation (grade buckets formed within training
folds only; Section~\ref{sec:robustness}) confirms the result with a slightly attenuated but still
significant increment ($z=3.16$, $p=1.6\times10^{-3}$).

\paragraph{The pooled increment is horizon-dependent and partly a pooling effect.} We stress that this
pooled significance must be read with care, because the five horizon files pool observations that are
not independent: the same firm can appear across files, so the effective sample is smaller than the
nominal 43{,}405, and the data carry no firm identifier with which to cluster or to form firm-grouped
cross-validation folds (a limitation we state in Section~\ref{sec:discussion}). Splitting the pooled
test into its five single-horizon files---each an independent cross-section---shows that the increment
is \emph{not} uniformly significant: at the three longer horizons the DeLong increment is
insignificant ($z=1.4$, $-0.6$, $1.6$), and it becomes significant only at the two shortest,
near-default horizons ($z=4.2$ and $4.7$), where the overlay is most valuable (Table~\ref{tab:horizon}).
The pooled $z$ therefore reflects both sample inflation and a genuine concentration of the signal close
to default; we do not interpret it as evidence of a large, uniform incremental effect. This horizon
gradient---stronger closer to default---is itself consistent with the overlay behaving as a
near-term deterioration signal.

\begin{table}[t]
\centering
\caption{Incremental discrimination by forecast horizon (single-horizon files, independent
cross-sections). The pooled significance concentrates at the near-default horizons.}
\label{tab:horizon}
\begin{tabular}{@{}lrrrc@{}}
\toprule
Horizon file & $n$ & $\Delta$AUC & DeLong $z$ & $p$ \\
\midrule
Pooled (5 horizons)        & 43{,}405 & $+0.0098$ & 7.77 & $7.8\times10^{-15}$ \\
\midrule
1 (5-year)                 & 7{,}027  & $+0.0045$ & 1.42 & 0.155 \\
2 (4-year)                 & 10{,}173 & $-0.0006$ & $-0.58$ & 0.563 \\
3 (3-year)                 & 10{,}503 & $+0.0035$ & 1.57 & 0.115 \\
4 (2-year)                 & 9{,}792  & $+0.0094$ & 4.24 & $2.2\times10^{-5}$ \\
5 (1-year)                 & 5{,}910  & $+0.0136$ & 4.66 & $3.1\times10^{-6}$ \\
\botrule
\end{tabular}
\end{table}

\paragraph{Economic significance.} Statistical significance aside, the \emph{economic} magnitude of the
pooled increment is modest, as expected when the base grade already exploits 62 ratios. Over the whole
sample the integrated discrimination improvement is small (IDI $=0.004$) and the Brier skill score is
0.4\%; the continuous net reclassification index is in fact slightly negative overall ($-0.06$), and the
minimum-misclassification-cost reduction at plausible cost ratios ($5$:$1$ to $20$:$1$) ranges only from
$-0.3\%$ to $+1.0\%$. The picture changes once attention is restricted to the middle band, where
interim re-rating is actually applied: the net reclassification index there is positive and
economically meaningful ($\text{NRI}_{\text{mid}}=+0.21$). The overlay is thus best read not as a
uniform accuracy gain but as a targeted reclassification tool whose value concentrates in the middle of
the spectrum---the same conclusion the band-moderation test reaches below.

\paragraph{The discrimination increment depends on discretizing the grade.} The base grade above is
discretized into seven buckets, which discards within-bucket variation in the underlying probability of
default and is the natural analogue of an agency letter scale. It is also, however, a modelling choice
that can create room for a continuous overlay to recover discarded signal. We therefore re-run the
nested test against a \emph{continuous} base probability of default---the base logit itself, with no
bucketing. Against this stronger benchmark the two-indicator overlay no longer improves ranking: the
out-of-fold AUC does not rise but instead falls, significantly so (pooled $\Delta\text{AUC}=-0.006$;
DeLong $z=-1.99$, $p=0.047$), while the likelihood-ratio statistic remains positive ($\text{LR}=247$). The
two readings are consistent once separated: the two indicators improve the \emph{calibration} of the
fitted default probability (hence the positive LR) but not its \emph{discrimination} once the grade is
continuous. We therefore do not claim a robust discrimination increment from the overlay; the AUC gains
in Table~\ref{tab:h1h2} are specific to a discretized grade and should be read as such. The findings we
do emphasize---indicator identity, asymmetry, and the band structure---do not rest on this increment,
as the robustness analysis shows.

\begin{table}[t]
\centering
\caption{Study 1: incremental discrimination over a discretized base grade (out-of-fold). Against a
continuous base PD the AUC increment vanishes; see text.}
\label{tab:h1h2}
\begin{tabular}{@{}lccc@{}}
\toprule
Model & AUC & Accuracy ratio & Brier \\
\midrule
$M_0$: base grade        & 0.729 & 0.458 & 0.0438 \\
$M_1$: base grade $+$ 2-proxy & 0.739 & 0.477 & 0.0436 \\
\midrule
\multicolumn{4}{@{}l}{DeLong $z=7.77$ ($p=7.8\times10^{-15}$); LR $=82.3$ ($p<10^{-18}$).}\\
\botrule
\end{tabular}
\end{table}

\begin{figure}[t]
\centering
\includegraphics[width=0.55\textwidth]{figures/fig04_study1_roc_m0_m1}
\caption{ROC curves for the base-grade model ($M_0$) and the augmented model ($M_1$).}
\label{fig:roc}
\end{figure}

\paragraph{A downside pattern in the point estimates (H3), reported as suggestive only.} Estimating
\eqref{eq:h3}, the point estimates are suggestive of asymmetry: deterioration relative to grade peers
enters strongly (coefficient $2.81$, $z=7.7$) while improvement is smaller and insignificant
(coefficient $1.22$, $z=1.5$). We are careful, however, not to overstate this. The correct test of
asymmetry is of \emph{equal slopes}, $\gamma^{+}-\gamma^{-}=0$, since under the hinge decomposition
$d=d^{+}+d^{-}$ a symmetric response is a single common slope; and that contrast is \emph{not}
significant ($\gamma^{+}-\gamma^{-}=1.6$, $p=0.11$). The improvement coefficient is same-signed and
non-negligible, so the honest statement is that improvement is estimated less precisely than
deterioration, not that it is zero. Two further facts counsel caution. The deviation $d^{+}$ has a
substantially wider spread than $d^{-}$ (standard deviations $0.052$ versus $0.031$), because the fitted
signal is right-skewed at a $4.8\%$ event rate, so a larger and better-identified downside slope can
arise from the distribution of $d$ alone, with no payoff asymmetry; and against a continuous base
probability of default the sign pattern reverses (improvement becomes the significant side), which is
the signature of a bucketing artifact rather than an economic effect. We therefore report the downside
pattern as a suggestive regularity---consistent in direction with the market evidence that bad credit
news moves prices more than good \citep{holthausen1986,dichev2001}---but do \emph{not} claim it as an
established asymmetry, and it is not part of the paper's contribution. Table~\ref{tab:h3robust} records
how the pattern behaves across specifications; the one specification in which the two slopes differ
non-trivially (the non-solvency placebo pair) does so in the \emph{opposite} direction and
insignificantly, underscoring that the pattern is not robust enough to build a claim on.

\begin{table}[t]
\centering
\caption{The downside point-estimate pattern (H3) across specifications, reported as suggestive. The
correct asymmetry test ($\gamma^{+}-\gamma^{-}=0$) is not rejected in the baseline ($p=0.11$); the
pattern reverses against a continuous base PD.}
\label{tab:h3robust}
\begin{tabular}{@{}lrrc@{}}
\toprule
Specification & $z(d^{+})$ & $z(d^{-})$ & $\gamma^{+}-\gamma^{-}=0$ \\
\midrule
Baseline (logit, discretized grade) & 7.7 & 1.5 & not rejected ($p=0.11$) \\
Probit (symmetric link)             & 8.7 & 0.5 & direction unchanged \\
Global-mean deviation               & 9.9 & $-0.2$ & direction unchanged \\
Non-solvency placebo pair           & 0.7 & $-1.7$ & reverses, insignificant \\
Continuous base PD                  & 1.6 & $-2.6$ & reverses \\
\botrule
\end{tabular}
\end{table}

\paragraph{Band moderation (H4).} The interaction of the signal with the middle-band indicator in
\eqref{eq:h4} is positive and significant (coefficient 5.86, $z=5.06$): the incremental value of the
two indicators is concentrated in the middle of the grade distribution. Because logit interaction
coefficients are not directly interpretable in magnitude, we also report the average marginal effect of
the overlay signal, which is about three times larger in the middle band than at the extremes---the
same conclusion on the probability scale. Figure~\ref{fig:h4} shows the incremental AUC is markedly
larger in the middle band than at the extremes. We note, however, that this interaction is contingent
on the quantile-bucketed grade (Section~\ref{sec:robustness}); the firmer evidence for the middle-band
pattern is the Study~2 recovery result (H5), which does not rely on the Study~1 bucketing.

\begin{figure}[t]
\centering
\includegraphics[width=0.5\textwidth]{figures/fig06_study1_h4_band_gain}
\caption{Incremental AUC ($M_1-M_0$) within the middle band versus the extremes.}
\label{fig:h4}
\end{figure}

\subsection{Study 2: rating recovery and information loss}
Table~\ref{tab:recovery} reports rating recovery. Two indicators recover a quadratic weighted kappa
of 0.345 against 0.654 for the full 25-ratio model, and within-one-grade accuracy of 77.0\% against
90.9\%---about half of the full information, echoing the ``intermediate zone'' of Study~1
(Figure~\ref{fig:qwk}). The two kappas are well separated: their 2{,}000-replication bootstrap 95\%
confidence intervals, $[0.625,0.680]$ for the full model and $[0.303,0.384]$ for the two indicators, do
not overlap. In the ordered logit \eqref{eq:ordlogit}, both proxies are significant with
the expected signs---leverage enters positively (debt ratio 0.52, $z=11.9$) and profitability
negatively (return on assets $-0.87$, $z=-3.3$)---and remain so under agency fixed effects.

\begin{table}[t]
\centering
\caption{Study 2: rating recovery (random forest, five-fold).}
\label{tab:recovery}
\begin{tabular}{@{}lccc@{}}
\toprule
Model & QWK & Accuracy & Within-1 grade \\
\midrule
Full (25 ratios)       & 0.654 & 0.509 & 0.909 \\
2 proxies (leverage, profitability) & 0.345 & 0.315 & 0.770 \\
\botrule
\end{tabular}
\end{table}

\begin{figure}[t]
\centering
\includegraphics[width=0.45\textwidth]{figures/fig07_study2_recovery_qwk}
\hfill
\includegraphics[width=0.5\textwidth]{figures/fig08_study2_band_within1}
\caption{Left: rating recovery (QWK), full versus two indicators. Right: within-one-grade recovery by band.}
\label{fig:qwk}
\end{figure}

\paragraph{Information loss by band (H5).} The within-one-grade information lost by using only two
indicators is larger at the extremes (18.4 percentage points) than in the middle band (11.6
percentage points). This band gap is statistically distinguishable from zero, not merely a raw
difference: a 2{,}000-replication bootstrap of the within-one-grade accuracy gap between the middle
band and the extremes gives a point estimate of about 13.7 percentage points with a 95\% confidence
interval of $[9.3,\,18.0]$, which excludes zero. The migration matrix (Figure~\ref{fig:migration})
shows the two-indicator model pulls firms toward the center and rarely predicts the extremes, which is
where its recovery is weakest. H5 is supported: the two indicators are least adequate precisely where
interim re-rating is, in practice, withheld.

\begin{figure}[t]
\centering
\includegraphics[width=0.6\textwidth]{figures/fig09_study2_migration_heatmap}
\caption{Grade migration: actual grade versus the grade predicted from two indicators.}
\label{fig:migration}
\end{figure}

%%=====================================================%%
\section{Robustness and model-risk analysis}\label{sec:robustness}

The analyses in this section double as a model-risk assessment of the overlay: they probe the three
specification choices on which a reduced-form overlay could fail---how the base grade is discretized
(discretization risk), which estimator maps indicators to grades (estimator risk), and which indicators
enter the signal (specification risk). Reporting all three is the kind of sensitivity and limitation
testing that model-validation standards expect \citep{bcbs2005validation,sr2011modelrisk}.

\subsection{Bucket-boundary sensitivity}
The base grade in Study~1 discretizes a continuous default probability, so we vary the number of
buckets (five, seven, ten) and the boundary rule (quantile versus fixed probability-of-default cuts).
Under quantile grading the incremental result is stable: the DeLong test is significant at all three
bucket counts and the H4 interaction remains positive with $z$ between 4.7 and 5.2
(Figure~\ref{fig:bucket}). Under arbitrary fixed cuts the base grade becomes weak (AUC 0.56) and the
overall incremental effect inflates, but the band definition is distorted and the H4 interaction flips
sign. We read this plainly: the asymmetry result (H3) is robust to bucket count, but the band-moderation
result (H4) is contingent on balanced, quantile-based grades and does not survive arbitrary cuts. We
therefore treat H4 as suggestive rather than firmly established, and note that the independent
middle-band evidence from Study~2 (H5), which does not rely on the Study~1 bucketing and is supported by
a bootstrap confidence interval, is the stronger basis for the band claim.

\begin{figure}[t]
\centering
\includegraphics[width=0.6\textwidth]{figures/fig10_robust_bucket_h4z}
\caption{H4 interaction $z$-statistic across grade-bucket schemes; dashed lines mark $\pm1.96$.}
\label{fig:bucket}
\end{figure}

\subsection{A regularized ordinal benchmark and cross-estimator agreement}
The high-dimensional maximum-likelihood ordered logit with full fixed effects fails to converge, so we
re-estimate the full rating model with an $L_2$-penalized proportional-odds specification
(Appendix~\ref{app:po}), which converges and preserves the ordering (full-model QWK 0.40 versus 0.13
for two indicators). Crucially, the band asymmetry underlying H5 replicates across three different
estimators---random forest, maximum-likelihood ordered logit, and the regularized proportional-odds
model---each showing larger two-indicator information loss at the extremes than in the middle
(Figure~\ref{fig:h5three}). The conclusion does not hinge on a single estimator.

\begin{figure}[t]
\centering
\includegraphics[width=0.62\textwidth]{figures/fig11_robust_h5_three_methods}
\caption{Two-indicator information loss (within-one-grade) by band across three estimators.}
\label{fig:h5three}
\end{figure}

\subsection{Leakage-free re-estimation and imbalance}
Because the base grade discretizes a fitted probability, we re-estimate Study~1 in a fully nested
scheme: within each training fold we fit the base model, form the quantile bucket edges, fit the
signal, and fit $M_0$ and $M_1$, then predict the held-out fold, so no test information enters the grade
or its buckets. The increment survives: the out-of-fold AUC still rises under the fully nested scheme
with a DeLong $z=3.16$ ($p=1.6\times10^{-3}$)---attenuated relative to the in-sample-bucket version but
still significant---confirming that the pooled result is not an artifact of bucket leakage. The default rate of
4.8\% is handled by class weighting; AUC and the accuracy ratio are stable across treatments (balanced
weighting 0.740/0.480 versus event oversampling 0.743/0.485), while an unweighted model is weaker
(0.671), so we retain balanced weighting throughout.

\subsection{Indicator choice}
We first replace the leverage--profitability pair with a liquidity--activity pair (current ratio and
asset turnover) in both studies. The contrast is stark (Table~\ref{tab:indicator},
Figure~\ref{fig:indicator}): in Study~1 the liquidity--activity signal adds nothing beyond the base
grade (likelihood-ratio statistic 0.7, $p=0.39$), against 82.3 for leverage--profitability; in Study~2
its rating-recovery QWK is 0.145 against 0.345.

To show this is not an artifact of a single hand-picked pair, we compute the incremental likelihood
ratio of every pair drawn from twelve ratios spanning leverage, profitability, liquidity, and
activity, over a fixed full-model grade (66 pairs; Figure~\ref{fig:pairs}). The ordering is wide and
systematic: the incremental likelihood ratio spans from near zero to about 90, so some pairs add two
orders of magnitude more information than others, even though every pair adds exactly two regressors.
We are deliberately careful about \emph{what} this ordering reflects, because it is easy to over-read.
It is not a clean ``solvency beats liquidity'' category effect: a permutation test on the broad
category labels does not reject the null that the solvency category and the liquidity/activity category
have the same mean incremental LR ($p=0.80$), and the highest-ranking pairs are in fact those containing
one particular high-information working-capital ratio, while some textbook leverage pairs rank near the
bottom. What the distribution shows is therefore an \emph{informativeness} ordering, not a category
ordering: the pairs that carry outcome-relevant information dominate, and the broad accounting label of
a ratio is a poor guide to whether it does. This is the right reading of the headline result. A sceptic
might object that \emph{any} two extra predictors will add a statistically detectable, if tiny,
increment beyond a fixed grade, so that the increment in Table~\ref{tab:h1h2} is mechanical. The
pairwise distribution refutes the strong form of this: the 66 pairs do not cluster at a small positive
value but span a wide range whose lower end is statistically indistinguishable from zero, functioning as
a within-design placebo. Two indicators add little when they carry little information and a great deal
when they carry much; the increment is driven by indicator \emph{identity}---which indicators enter---not
by the act of adding two regressors. Of the 66 tests, 56 are significant at the 5\% level without
correction, 51 survive a conservative Bonferroni threshold, and 56 survive a Benjamini--Hochberg
false-discovery-rate control, so the spread of the distribution is not an artifact of multiplicity. The
one caveat we flag is that this pairwise distribution is estimated on the pooled Polish panel, which is
not firm-independent (Section~\ref{sec:discussion}); the external replication of
Section~\ref{sec:multidata} provides the firm-independent evidence that the identity-over-count pattern
is not a pooling artifact. In corporate-default settings the informative ratios happen to be solvency
ratios, the dimensions emphasized since \citet{altman1968}, but we do not claim solvency is privileged
as a \emph{category}---only that informativeness, of which solvency is the default-relevant instance,
governs the overlay. The downside pattern, by contrast, holds for neither pair strongly enough to be
established (H3), so we do not rest anything on it.
to indicator choice.

\begin{figure}[t]
\centering
\includegraphics[width=0.85\textwidth]{figures/fig13_review_all_pairs_LR}
\caption{Incremental likelihood ratio of each two-ratio signal over a fixed full-model grade (66 pairs
from twelve ratios), ordered by informativeness. The ordering reflects the information a pair carries,
not its broad accounting category (a permutation test on the solvency versus liquidity/activity label
gives $p=0.80$); the pair used in the main analysis is shaded darker. Horizontal bars mark medians.}
\label{fig:pairs}
\end{figure}

\begin{table}[t]
\centering
\caption{Indicator-choice comparison across both studies.}
\label{tab:indicator}
\begin{tabular}{@{}lcc@{}}
\toprule
Indicator pair & Study 1 incremental LR & Study 2 recovery QWK \\
\midrule
Leverage $+$ profitability & 82.3 & 0.345 \\
Liquidity $+$ activity     & 0.7  & 0.145 \\
\botrule
\end{tabular}
\end{table}

\begin{figure}[t]
\centering
\includegraphics[width=0.7\textwidth]{figures/fig12_robust_indicator_choice}
\caption{Indicator choice: Study 1 incremental LR (left) and Study 2 recovery QWK (right).}
\label{fig:indicator}
\end{figure}

\subsection{External replication on independent cross-sections}\label{sec:multidata}
The Polish panel's one limitation we cannot repair internally---pooled, repeated-firm observations with
no firm identifier---motivates a external check on data that are free of it. We therefore replicate the
two central tests on four additional, firm-independent cross-sections from the consumer-credit and
classification literature: a peer-to-peer consumer-loan dataset ($n=32{,}581$), the ``Give Me Some
Credit'' delinquency data ($n=150{,}000$), the Taiwan credit-card default data ($n=30{,}000$), and, as a
non-credit control, the Adult income-classification data ($n=48{,}842$). Each is a single cross-section
of independent records, so the pooling concern of Section~\ref{sec:results} does not arise. On each we
build a base grade from all other features, discretize it into seven buckets exactly as in
Study~1, and test (i) whether a two-indicator overlay adds out-of-fold discrimination, and (ii) whether
a chosen two-feature pair outperforms random two-feature placebo pairs.

One pattern replicates cleanly and one is more nuanced, and we report both rather than only the
favourable reading (Table~\ref{tab:multidata}). First, a two-indicator overlay adds significant
out-of-fold discrimination in every independent dataset (DeLong $z$ from 11.8 to 44.1), including the
non-credit Adult control, so the incremental-information finding is not special to the Polish sample---but,
precisely because it also holds in the non-credit control, it confirms that ``two informative features
add something to a discretized grade'' is a \emph{generic} property, not evidence for a credit-specific
mechanism. Second, the placebo contrast must be read carefully. The chosen pair beats the median random
pair by large factors in three of four datasets, but the relevant content differs: in the Polish panel
and the rating cross-section the dominant pair is genuinely a solvency pair (leverage and profitability),
whereas in the Taiwan data the dominant pair pairs the credit limit with the most recent repayment-delay
status---a near-outcome variable, not a solvency ratio---so its dominance reflects proximity to the
outcome rather than solvency content. In the Give Me Some Credit data the solvency-type pair
(debt ratio and income) in fact \emph{loses} to the median random pair (ratio 0.3). The honest,
cross-dataset generalization is therefore the weaker one: within a two-feature overlay, the
\emph{identity} of the pair---whether it carries the outcome-relevant information---dominates the mere
fact of using two features, and in genuine corporate-default settings (Polish, ratings) that
outcome-relevant information is solvency. We do not claim that solvency content universally governs the
increment; the consumer datasets show the operative property is informativeness, of which solvency is
the default-relevant instance.

\begin{table}[t]
\centering
\caption{External replication on four independent cross-sections (plus the single-horizon Polish file).
A two-indicator overlay adds out-of-fold discrimination in every independent dataset. The final column
is the incremental-LR ratio of the chosen pair to the median random placebo pair; the chosen pair and
whether it is a solvency pair are noted, since the dominant pair is a solvency pair only in the
genuine-default settings.}
\label{tab:multidata}
\small
\begin{tabular}{@{}lrrrcl@{}}
\toprule
Dataset & $n$ & $\Delta$AUC & $z$ & LR ratio & Chosen pair (solvency?) \\
\midrule
Consumer loans (P2P)       & 32{,}581  & $+0.092$ & 29.2 & 21.4 & debt burden $+$ income \\
Give Me Some Credit         & 150{,}000 & $+0.008$ & 11.8 & 0.3  & debt ratio $+$ income (yes) \\
Taiwan credit card          & 30{,}000  & $+0.041$ & 13.8 & 24.8 & limit $+$ recent delay (no) \\
Adult (control)             & 48{,}842  & $+0.088$ & 44.1 & 4.7  & education $+$ hours (non-credit) \\
\midrule
Polish (1-yr, single)       & 7{,}027   & $+0.005$ & 1.42 & 7.9 & leverage $+$ profitability (yes) \\
\botrule
\end{tabular}
\end{table}

%%=====================================================%%
\section{Discussion}\label{sec:discussion}

Our results speak to the design of interim re-rating policies. The central message is that a
parsimonious point-in-time overlay is informative, but conditionally so, and the conditions are
economically interpretable.

\paragraph{Band targeting is justified.} The incremental value of the two-indicator signal is
concentrated in the middle band (H4), and the rating information that two indicators fail to recover is
largest at the extremes (H5). Both results, obtained on different datasets with different outcomes,
point the same way: confining interim re-rating to the middle of the rating spectrum---common practice
---loses little, because near-default and high-grade firms are precisely where two indicators are least
adequate. The practice that is often defended on risk-aversion grounds has a measurement rationale.

\paragraph{A possible downside tilt, not established.} The point estimates hint that deterioration is
more informative than improvement, which would mirror the long-standing market evidence that downgrades
move prices more than upgrades \citep{holthausen1986,dichev2001,norden2004}. We are careful not to build
on this: as Section~\ref{sec:results} shows, the correct equal-slopes test does not establish the
asymmetry, the downside deviation is more dispersed than the upside, and the pattern reverses against a
continuous base probability of default. We therefore do not recommend an asymmetric, downside-only use
of the overlay on the strength of our evidence; whether such a tilt is real is a question for data with
a genuine two-period structure. What our evidence does support about \emph{where} the overlay informs is
the band pattern of the previous paragraph, which rests on the independent Study~2 recovery result.

\begin{figure}[t]
\centering
\includegraphics[width=0.95\textwidth]{figures/fig_band_schematic}
\caption{Where the overlay informs. Its incremental value concentrates in the middle band
(BBB--B), where the overlay is applied, and the two-indicator information loss is largest at the rating
extremes, where it is withheld. The direction of the effect (whether deterioration is more informative
than improvement) is not statistically established (H3) and is therefore not depicted.}
\label{fig:band}
\end{figure}

\paragraph{Indicator selection is a first-order choice---and the main model risk.} The overlay works
with leverage and profitability but not with liquidity and activity. Because an interim re-rating scheme
is defined as much by its two indicators as by its existence, the selection of those indicators deserves
the same scrutiny as the decision to re-rate at all. This connects to the feature-selection literature
\citep{tsai2009,hajek2013,laborda2021}: in a radically reduced model, identity dominates count. In
model-risk terms, the choice of indicators---not the discretization or the estimator, both of which leave
the qualitative results intact---is where the overlay is most exposed to specification error, so it is the
choice a validation process should scrutinize most closely \citep{danielsson2002,sr2011modelrisk}. The
pairwise analysis sharpens this further: it is specifically the \emph{solvency} content (leverage, then
profitability) that carries the overlay, so an interim scheme should be built around solvency indicators
rather than liquidity or turnover.

\paragraph{Implications, stated as conditional conjectures.} The implications below follow from the
measured \emph{shape} of the information---where and in which direction two indicators inform---rather
than from an observed re-rating process, and we state them as conjectures to be tested on data with a
genuine two-period structure, not as validated policy prescriptions. Subject to that caveat, three
implications suggest themselves. First, because the overlay's incremental value is real but
economically modest in aggregate and concentrated in the middle band and on the downside, a reduced
two-indicator signal looks better suited to flagging firms for fuller review when two solvency
indicators deteriorate than to symmetric, whole-spectrum re-grading. Second, the band result suggests
that such a signal adds little for investment-grade and near-default names, where two indicators are
least adequate, so a reduced overlay should not drive mechanical actions there. Third, because the
increment rests on indicator identity rather than count, a reduced scheme should fix and validate its
solvency indicators explicitly rather than treat a two-indicator model as safe merely because it is
small. Whether these conjectures carry over to SME and supply-chain settings, in which buyers re-score
many suppliers cheaply and frequently \citep{altman2005emerging,berger2006,giannetti2011}, is an
empirical question for data from those settings; our evidence speaks to the shape of the information,
not to its magnitude in any particular market.

\paragraph{Limitations.} Five limitations bound the interpretation. First and most important, the
public datasets do not contain an observed annual-versus-interim timeline, so the two-clock structure is
a \emph{reconstruction}: we operationalize it by separating the indicator set and, in robustness, by the
forecast-horizon gradient, but we cannot observe a single firm's own mid-year revision. The construct is
therefore an interpretation of a feature-parsimony result, and confirming it on data with a genuine
two-period structure is the first-order task for future work. Second, the pooled Polish files do not
carry a firm identifier, so the same firm may recur across the five horizon files and the pooled
observations are not independent. We cannot cluster standard errors by firm or form firm-grouped
cross-validation folds on these data; the pooled significance should therefore be read as an upper
bound, and we rely instead on the single-horizon decomposition (Table~\ref{tab:horizon}) and on the
four firm-independent cross-sections of Section~\ref{sec:multidata}, where no such dependence exists, to
establish that the qualitative findings do not hinge on the pooling. Third, the Study~1 base grade is a
constructed probability-of-default grade rather than an agency grade, so its boundaries are a modeling
choice (we show the core results are robust to them and are leakage-free). Fourth, Study~2 has agency
grades but no default outcome, so its evidence concerns information recovery, not predictive accuracy.
The two datasets share no firms and overlap in neither country nor period, so they corroborate rather
than replicate each other: agreement between them is mutual support across independent samples and
outcomes, not a re-estimate of one quantity, and we read it only that way. Fifth, the data are Polish firms and a
multi-agency, largely US rating set rather than a domestic SME or single-market setting, so the external
validity of the specific magnitudes---though not, we argue, of the qualitative shape---awaits replication
on SME and supply-chain data; our claims are methodological, about the \emph{shape} of the information
rather than its exact level.

%%=====================================================%%
\section{Conclusion}\label{sec:conclusion}

We asked, for a radically reduced credit overlay of only two indicators, whether its specification
risk comes from how many indicators it uses or from which, using public datasets and a fully
reproducible, out-of-fold design. Our central finding is that identity dominates count. Across 66
indicator pairs the incremental likelihood ratio spans two orders of magnitude; the pairs that carry
outcome-relevant information sit at the top and the uninformative ones are indistinguishable from zero,
and the ordering survives Bonferroni and false-discovery control. The ordering reflects informativeness
rather than a broad accounting category---a permutation test does not find solvency ratios to outrank
liquidity ratios as categories---so it is which indicators carry information, not their textbook label
or their number, that matters; in corporate-default settings the informative ones are solvency ratios.
Pairs that carry information behave entirely differently from those that do not, though both add exactly
two regressors, so in a reduced-form overlay indicator selection---not indicator count---is the
first-order source of specification risk. Two further findings frame this one. The overlay's incremental ranking power is
modest and fragile: significant on the pooled panel but concentrated near default, and not robust to a
continuous base probability of default, so we claim no uniform discrimination gain. And a nonparametric
rating-recovery analysis shows two indicators are least adequate at the rating extremes and most
adequate in the middle band, a gap backed by a bootstrap confidence interval. A downside-asymmetry
pattern appears in the point estimates but, once the correct equal-slopes contrast is tested, is not
statistically established, and we report it only as suggestive.

Taken together, the results characterize what governs the information in a radically reduced,
two-indicator overlay: its value turns on indicator identity, and its recovery is weakest at the rating
extremes. As a conditional conjecture for practice---to be tested on data with a genuine two-period
structure---this is consistent with building a reduced overlay around middle-band firms and solvency
indicators, and with validating the \emph{choice} of those indicators as the point of greatest
specification risk rather than treating a two-indicator model as safe merely because it is small. For
research, the paper offers a public-data template for testing which indicators, not how many, carry the
information in parsimonious credit overlays without proprietary agency data. The natural next step is to replicate the design on data with an observed two-period timeline,
which would let the magnitudes---not only the shape---of the information be estimated directly.

%%=====================================================%%
\backmatter

%% Acknowledgements omitted in the double-blind version (add on acceptance).

\section*{Declarations}

\noindent\textbf{Funding.} The authors received no specific funding for this work.

\medskip
\noindent\textbf{Competing interests.} The authors declare no competing financial or non-financial
interests.

\medskip
\noindent\textbf{Data availability.} Both datasets are public: the Polish companies bankruptcy data
\citep{zieba2016} distributed through the UCI Machine Learning Repository, and the corporate
credit-rating data (OpenML dataset \#43344). All code, notebooks, and figures needed to reproduce every
result and exhibit are available at an anonymized repository,
\url{https://anonymous.4open.science/r/parsimonious-interim-rerating-6B71}; the repository documents
how to obtain the data, and analyses run on fixed random seeds, so the reported figures are
deterministic.

\medskip
\noindent\textbf{Declaration of generative AI and AI-assisted technologies.} During the preparation of
this manuscript, the author used Claude (Anthropic), a generative AI assistant, to help draft and edit
the text, to write and run the Python analysis code, to prepare the figures, and to search and organize
the cited literature. After using this tool, the author reviewed, verified, and edited all
content---including the study design, the data analysis and numerical results, the figures, and the
references---and takes full responsibility for the content of the publication.

%%=====================================================%%
\begin{appendices}

\section{Proxy indicator definitions}\label{app:indicators}

Table~\ref{tab:defs} lists the two proxy indicators used in each study and their robustness
counterparts. In Study~1 the ratios are items of the 64-ratio set of \citet{zieba2016}; in Study~2 they
are items of the 25-ratio set in the rating data.

\begin{table}[h]
\centering
\caption{Proxy and alternative indicators.}
\label{tab:defs}
\begin{tabular}{@{}lll@{}}
\toprule
Role & Study 1 (Polish) & Study 2 (ratings) \\
\midrule
Profitability (main) & net profit / total assets & return on assets \\
Leverage (main)      & total liabilities / total assets & debt ratio \\
Liquidity (alt.)     & current assets / short-term liabilities & current ratio \\
Activity (alt.)      & sales / total assets & asset turnover \\
\botrule
\end{tabular}
\end{table}

\section{$L_2$-penalized proportional-odds model}\label{app:po}

To benchmark the full rating model without the non-convergence of a high-dimensional
maximum-likelihood ordered logit, we estimate a cumulative-logit (proportional-odds) model with an
$L_2$ penalty on the slope coefficients. For $K$ ordered grades the cumulative probabilities are
\begin{equation}
\Pr(R_i \le k) = \Lambda(\tau_k - \boldsymbol{\beta}^{\top}Z_i), \qquad k = 1,\dots,K-1,
\end{equation}
with $\Lambda$ the logistic function and ordered cutpoints $\tau_1 < \cdots < \tau_{K-1}$ enforced by a
cumulative-exponential parameterization. The penalized negative log-likelihood is
\begin{equation}
\ell_{\lambda}(\boldsymbol{\beta},\boldsymbol{\tau}) = -\sum_{i} \log\Pr(R_i = r_i) + \lambda\,\lVert\boldsymbol{\beta}\rVert_2^2,
\end{equation}
minimized by L-BFGS with $\lambda=2$. This specification converges with the full 25-ratio set plus
agency and sector fixed effects and yields a full-model QWK of 0.40 against 0.13 for the two-indicator
model, with band-wise information loss of 4.2 percentage points in the middle band and 8.3 at the
extremes---the same asymmetry found with the random forest and the maximum-likelihood ordered logit
(Section~\ref{sec:robustness}). To validate the implementation, we compare it on the two-indicator case
(where a standard maximum-likelihood ordered logit converges) against the \texttt{statsmodels}
\texttt{OrderedModel}: the slope coefficients agree to within $10^{-4}$ (debt ratio $0.473$ vs $0.473$;
return on assets $-0.612$ vs $-0.612$), with identical signs and ordering.

\section{Additional tables}\label{app:tables}

Table~\ref{tab:gradient} reports the out-of-fold AUC of the base-grade and augmented models by
forecast horizon (Study~1). The incremental gain is retained as the horizon shortens, and is largest at
the one-year-ahead horizon, consistent with an interim signal being more valuable closer to default.

\begin{table}[h]
\centering
\caption{Study 1: out-of-fold AUC by forecast horizon.}
\label{tab:gradient}
\begin{tabular}{@{}lccc@{}}
\toprule
Horizon & $M_0$ (base) & $M_1$ ($+$2-proxy) & Gain \\
\midrule
1 (5-year) & 0.718 & 0.727 & 0.009 \\
2          & 0.662 & 0.669 & 0.007 \\
3          & 0.718 & 0.729 & 0.011 \\
4          & 0.730 & 0.740 & 0.010 \\
5 (1-year) & 0.810 & 0.822 & 0.011 \\
\botrule
\end{tabular}
\end{table}

\end{appendices}

%%=====================================================%%
\bibliography{references}

\end{document}
